\section{相关工作}
\label{sec:related-work}

\paragraph{长期记忆基准。}
LoCoMo 和 LongMemEval 评估长对话、时序更新和多会话检索 \citep{maharana2024locomo,wu2025longmemeval}。MEME 增加删除、级联影响和缺失事实问题 \citep{jung2026meme}。MemoryAgentBench（MAB）在一个评分框架中同时测试准确检索、时序与测试时学习、长程理解和选择性遗忘 \citep{hu2025memoryagentbench}。这些基准分别覆盖存储和读取的不同部分，因此本文将它们作为互补测试，而不把长期记忆压缩成一个分数。

\paragraph{智能体记忆系统。}
MemGPT 将上下文管理描述为工作窗口与外部存储之间的分页 \citep{packer2023memgpt}。Model-Native Computing Architecture 进一步把上下文容量、缓存复用和智能体调度放在系统级资源管理框架下讨论 \citep{lin2026icam}。Mem0、Zep/Graphiti、HippoRAG-2、A-MEM、MIRIX 和 Cognee 使用摘要、笔记、图或它们的组合来保存跨会话信息 \citep{chhikara2025mem0,rasmussen2025zep,gutierrez2025hipporag2,xu2025amem,wang2025mirix,markovic2025cognee}。这些系统在“什么被保存为记忆”上各有不同：Deep-Dream 保留源文档，并把概念和关系作为指向源文档的链接，因此图中的派生节点是检索入口，而不是替代事实。

\paragraph{时序与版本化知识。}
时序知识图谱表示不断变化的实体和关系；概率记录链接则区分错误合并与漏匹配 \citep{fellegi1969theory}。Deep-Dream 将这些思想用于智能体记忆：处理时间排列观测和版本，延迟对齐使候选身份在证据充分之前保持可修订，源文本则始终保持不变。事件日期仍然属于观测内容，可以参与检索，但不会另建一套状态时钟。

\paragraph{压缩、检索与证据。}
信息瓶颈和率失真理论研究紧凑表示与目标信息之间的取舍 \citep{tishby2000information,shannon1959coding}。联想记忆研究如何从部分线索中检索内容 \citep{ramsauer2021hopfield}，检索增强生成则从外部来源选择生成证据 \citep{lewis2020rag}。长上下文模型在相关片段埋藏于大窗口中时可能丢失信息 \citep{liu2024lost}；归因研究检查回答是否引用了来源 \citep{gao2023alce}。Deep-Dream 将这些思想放入持久化记忆：概念和关系缩小搜索范围，源文本保留原始措辞，智能体决定读取多少内容；第~\cref{sec:information-view} 节给出这一分工的形式化分析。

\paragraph{智能体检索。}
A-RAG 将关键词、语义和片段读取工具交给推理模型，用于多跳检索 \citep{du2026arag}。Deep-Dream 同样允许模型选择下一步读取动作，但面向的是不断变化的个人历史。它进一步需要处理身份修订、版本链接、跨文档 Episode 和基于原文的验证。
